\documentclass[a4paper,12pt]{article}

\usepackage[french]{babel}
\usepackage[T1]{fontenc}
\usepackage[utf8]{inputenc}
\usepackage[top=2.5cm, bottom=2.5cm, left=2.5cm, right=2.5cm]{geometry}
\usepackage{graphicx}
\usepackage{listings}
\usepackage{xcolor}
\usepackage{fancyhdr}
\usepackage{hyperref}
\usepackage{float}
\usepackage{booktabs}
\usepackage{mdframed}
\usepackage{parskip}
\usepackage{titlesec}

% ---- Couleurs ----
\definecolor{javagreen}{rgb}{0.25,0.5,0.35}
\definecolor{javapurple}{rgb}{0.5,0,0.35}
\definecolor{javaocre}{rgb}{0.65,0.4,0.0}
\definecolor{codebg}{rgb}{0.97,0.97,0.97}
\definecolor{rouge}{RGB}{220,53,69}
\definecolor{vert}{RGB}{40,167,69}
\definecolor{gris}{RGB}{180,180,180}

% ---- Style Java ----
\lstdefinestyle{java}{
  language=Java,
  backgroundcolor=\color{codebg},
  basicstyle=\ttfamily\footnotesize,
  keywordstyle=\color{javapurple}\bfseries,
  commentstyle=\color{javagreen}\itshape,
  stringstyle=\color{javaocre},
  numberstyle=\tiny\color{gray},
  numbers=left,
  stepnumber=1,
  frame=single,
  rulecolor=\color{gray!40},
  breaklines=true,
  showstringspaces=false,
  tabsize=4,
  captionpos=b,
  xleftmargin=1em,
}

% ---- Encadré rouge (test en échec) ----
\newcommand{\etatrouge}[1]{%
  \begin{mdframed}[linecolor=rouge,linewidth=2pt,backgroundcolor=rouge!8]%
    \textcolor{rouge}{\textbf{\textsc{Rouge} ---}} #1%
  \end{mdframed}%
}

% ---- Encadré vert (test qui passe) ----
\newcommand{\etatvert}[1]{%
  \begin{mdframed}[linecolor=vert,linewidth=2pt,backgroundcolor=vert!8]%
    \textcolor{vert}{\textbf{\textsc{Vert} ---}} #1%
  \end{mdframed}%
}

% ---- Screenshot (avec fallback si fichier absent) ----
\newcommand{\screenshot}[2]{%
  \begin{figure}[H]%
    \centering%
    \IfFileExists{screenshots/#1}{%
      \includegraphics[width=0.92\textwidth]{screenshots/#1}%
    }{%
      \fcolorbox{gris}{gris!20}{\parbox{0.88\textwidth}{%
        \centering\vspace{0.7cm}%
        \textcolor{gray}{\itshape [Screenshot à insérer : \texttt{#1}]}%
        \vspace{0.7cm}%
      }}%
    }%
    \caption{#2}%
  \end{figure}%
}

% ---- En-tête / pied de page ----
\pagestyle{fancy}
\fancyhf{}
\lhead{\small Qualité de développement — TP2}
\rhead{\small BUT Informatique 2\up{e} année}
\rfoot{\thepage}
\renewcommand{\headrulewidth}{0.4pt}

\hypersetup{colorlinks=true,linkcolor=black,urlcolor=blue}

% =====================================================================
\begin{document}
% =====================================================================

% ---- Page de titre ----
\begin{titlepage}
  \centering
  \vspace*{1.5cm}
  {\Large\bfseries Département Informatique --- BUT 2\up{e} année\par}
  \vspace{0.4cm}
  {\large Qualité de développement\par}
  \vspace{2cm}
  \rule{\textwidth}{1.5pt}\\[0.5cm]
  {\Huge\bfseries Compte rendu --- TP2\par}
  \vspace{0.3cm}
  {\Large Application de calcul statistique\\ pour groupes de TP\par}
  \rule{\textwidth}{1.5pt}\\[2cm]
  \begin{tabular}{ll}
    \textbf{Étudiants} & Ahmed Errebache \& Léa Couturas \\
    \textbf{Groupe}    & BUT Informatique 2A             \\
    \textbf{Date}      & 13 mai 2026                     \\
    \textbf{Dépôt}     & \texttt{QD3\_TP2\_AE\_LC}       \\
  \end{tabular}
  \vfill
  {\large IUT --- 2025--2026\par}
\end{titlepage}

\tableofcontents
\newpage

% =====================================================================
\section{Introduction}
% =====================================================================

\subsection{Contexte}

Dans le cadre du module \textit{Qualité de développement}, nous avons
développé une application Java permettant à un enseignant de saisir les
notes de son groupe TP, d'en calculer la moyenne et la médiane, puis de
stocker ce résultat avec ses identifiants.

Le sujet s'inspire du cas réel d'Alain Dupont (\texttt{alain.dupont@iut.fr}),
enseignant en informatique, qui avait jusqu'ici un simple tableau d'entiers
fixe. Notre objectif a été de concevoir une application propre, testée et
évolutive.

\subsection{Architecture du projet}

Le projet est structuré en trois classes métier séparées :

\begin{center}
\begin{tabular}{lp{9.5cm}}
  \toprule
  \textbf{Classe}         & \textbf{Rôle} \\
  \midrule
  \texttt{UtilisateurTab} & Représente l'enseignant : prénom, nom, adresse mail
                            validée par expression régulière, date de l'examen \\
  \texttt{CalculTab}      & Gère la liste des notes (\texttt{ArrayList<Integer>}),
                            calcule la moyenne, la médiane (avec tri interne), et
                            formate le résultat à stocker \\
  \texttt{Principale}     & Point d'entrée (\texttt{main}) : saisie interactive et
                            affichage des résultats \\
  \bottomrule
\end{tabular}
\end{center}

\subsection{Démarche TDD}

Nous avons adopté une approche \textbf{TDD} (Test-Driven Development) :

\begin{enumerate}
  \item Écriture du test --- il est \textbf{rouge} car le code n'existe pas encore.
  \item Implémentation minimale du code pour faire passer le test.
  \item Vérification --- le test passe, il est \textbf{vert}.
\end{enumerate}

Tous les tests suivent la convention de nommage
\texttt{givenÉtat\_whenAction\_thenComportement()}, sont annotés
\texttt{@DisplayName}, et utilisent \textbf{AssertJ} pour des assertions
lisibles.

% =====================================================================
\section{Modélisation de l'enseignant (\texttt{UtilisateurTab})}
% =====================================================================

\subsection{Objectif}

La classe \texttt{UtilisateurTab} représente l'enseignant. Une adresse mail
mal formatée doit être immédiatement rejetée par une
\texttt{IllegalArgumentException}, afin que le programme ne puisse jamais
stocker un résultat avec des identifiants invalides.

\subsection{Tests écrits}

\begin{lstlisting}[style=java, caption={\texttt{UtilisateurTabTest.java} --- extrait}]
@Test
@DisplayName("Creation reussie avec un email valide")
void constructor_createsUser_whenEmailIsValid() {
    // Given + When
    UtilisateurTab u = new UtilisateurTab(
            "Alain", "Dupont", "alain.dupont@iut.fr", DATE_EXAMEN);
    // Then
    assertThat(u.getPrenom()).isEqualTo("Alain");
    assertThat(u.getEmail()).isEqualTo("alain.dupont@iut.fr");
}

@ParameterizedTest(name = "Email invalide : \"{0}\"")
@DisplayName("Construction echoue pour des emails invalides")
@ValueSource(strings = { "pasunmail", "sans@domaine", "@nodomain.fr", "" })
void constructor_throwsException_whenEmailIsInvalid(String emailInvalide) {
    assertThatThrownBy(() ->
            new UtilisateurTab("Alain", "Dupont", emailInvalide, DATE_EXAMEN))
            .isInstanceOf(IllegalArgumentException.class)
            .hasMessageContaining("invalide");
}
\end{lstlisting}

\subsection{État rouge --- avant implémentation}

\etatrouge{4 tests échouent : aucune validation n'est encore présente,
\texttt{ajouterNote()} accepte n'importe quelle chaîne.}

\screenshot{AE_01_email_rouge.png}{%
  \texttt{mvn test -Dtest=UtilisateurTabTest} --- avant la validation email}

\subsection{Implémentation}

Nous avons ajouté une expression régulière compilée en constante statique
pour valider le format de l'email dès la construction de l'objet.

\begin{lstlisting}[style=java, caption={\texttt{UtilisateurTab.java} --- validation email}]
private static final Pattern EMAIL_PATTERN =
        Pattern.compile("[\\w.-]+@[\\w.-]+\\.[a-z]{2,}");

public UtilisateurTab(String prenom, String nom,
                      String email, LocalDate dateExamen) {
    if (!EMAIL_PATTERN.matcher(email).matches()) {
        throw new IllegalArgumentException(
                "Adresse mail invalide : " + email);
    }
    this.prenom = prenom;
    this.nom    = nom;
    this.email  = email;
    this.dateExamen = dateExamen;
}
\end{lstlisting}

\subsection{État vert --- après implémentation}

\etatvert{7 tests passent : création valide réussie et rejet des emails
incorrects (4 cas paramétrés).}

\screenshot{AE_01_email_vert.png}{%
  \texttt{mvn test -Dtest=UtilisateurTabTest} --- après la validation email}

% =====================================================================
\section{Calcul des statistiques (\texttt{CalculTab})}
% =====================================================================

% ----------------------------
\subsection{Calcul de la moyenne}
% ----------------------------

\subsubsection*{Objectif}

La méthode \texttt{calculerMoyenne()} calcule la somme de toutes les notes
divisée par leur nombre. Si la liste est vide, elle retourne \texttt{0.0}.

\subsubsection*{Tests écrits}

\begin{lstlisting}[style=java, caption={\texttt{CalculTabTest.java} --- tests moyenne}]
@Test
@DisplayName("Moyenne correcte pour plusieurs notes")
void calculerMoyenne_returnsCorrectAverage_forMultipleNotes() {
    // Arrange
    calcul.ajouterNote(10);
    calcul.ajouterNote(12);
    calcul.ajouterNote(14);
    // Act + Assert
    assertThat(calcul.calculerMoyenne()).isEqualTo(12.0);
}

@Test
@DisplayName("Moyenne vaut 0 quand aucune note n'est presente")
void calculerMoyenne_returnsZero_whenNoNote() {
    assertThat(calcul.calculerMoyenne()).isEqualTo(0.0);
}
\end{lstlisting}

\subsubsection*{État rouge --- avant implémentation}

\etatrouge{2 tests échouent : \texttt{calculerMoyenne()} retourne toujours
\texttt{0.0} quelle que soit la liste.}

\screenshot{AE_02_moyenne_rouge.png}{%
  \texttt{mvn test -Dtest=CalculTabTest} --- avant \texttt{calculerMoyenne()}}

\subsubsection*{Implémentation}

\begin{lstlisting}[style=java, caption={\texttt{CalculTab.java} --- calculerMoyenne()}]
public double calculerMoyenne() {
    if (notes.isEmpty()) return 0.0;
    int somme = 0;
    for (int note : notes) somme += note;
    return (double) somme / notes.size();
}
\end{lstlisting}

\subsubsection*{État vert --- après implémentation}

\etatvert{Tests de moyenne passent.}

\screenshot{AE_02_moyenne_vert.png}{%
  \texttt{mvn test -Dtest=CalculTabTest} --- après \texttt{calculerMoyenne()}}

% ----------------------------
\subsection{Calcul de la médiane}
% ----------------------------

\subsubsection*{Objectif}

La médiane est calculée sur une \textit{copie triée} de la liste, afin de
ne pas modifier l'ordre de saisie. Pour un nombre pair de notes, on fait
la moyenne des deux valeurs centrales.

\subsubsection*{Tests écrits}

\begin{lstlisting}[style=java, caption={\texttt{CalculTabTest.java} --- tests médiane}]
@Test
@DisplayName("Mediane : valeur centrale pour nombre impair de notes")
void calculerMediane_returnsMiddleValue_forOddCount() {
    // notes saisies en desordre intentionnellement
    calcul.ajouterNote(15); calcul.ajouterNote(10); calcul.ajouterNote(12);
    // [10, 12, 15] apres tri -> mediane = 12
    assertThat(calcul.calculerMediane()).isEqualTo(12.0);
}

@Test
@DisplayName("Mediane : moyenne des deux centraux pour nombre pair")
void calculerMediane_returnsAverageOfTwoMiddle_forEvenCount() {
    calcul.ajouterNote(10); calcul.ajouterNote(12);
    calcul.ajouterNote(14); calcul.ajouterNote(16);
    // (12 + 14) / 2 = 13.0
    assertThat(calcul.calculerMediane()).isEqualTo(13.0);
}
\end{lstlisting}

\subsubsection*{État rouge --- avant implémentation}

\etatrouge{3 tests échouent : \texttt{calculerMediane()} retourne toujours
\texttt{0.0}.}

\screenshot{AE_03_mediane_rouge.png}{%
  \texttt{mvn test -Dtest=CalculTabTest} --- avant \texttt{calculerMediane()}}

\subsubsection*{Implémentation}

\begin{lstlisting}[style=java, caption={\texttt{CalculTab.java} --- calculerMediane()}]
public double calculerMediane() {
    if (notes.isEmpty()) return 0.0;
    List<Integer> triees = new ArrayList<>(notes); // copie : ordre original conserve
    Collections.sort(triees);
    int milieu = triees.size() / 2;
    if (triees.size() % 2 == 1) return triees.get(milieu);
    return (triees.get(milieu - 1) + triees.get(milieu)) / 2.0;
}
\end{lstlisting}

\subsubsection*{État vert --- après implémentation}

\etatvert{10 tests passent dans \texttt{CalculTabTest}.}

\screenshot{AE_03_mediane_vert.png}{%
  \texttt{mvn test -Dtest=CalculTabTest} --- après \texttt{calculerMediane()}}

% =====================================================================
\section{Scénarios d'utilisation complète}
% =====================================================================

\subsection{Objectif}

Cinq scénarios \textit{end-to-end} valident le parcours complet d'Alain
Dupont : saisie d'un groupe de notes, calcul des statistiques et génération
de la chaîne à stocker. Ils utilisent la structure \texttt{@Nested} de JUnit 5.

\begin{center}
\begin{tabular}{clp{7cm}}
  \toprule
  \textbf{N°} & \textbf{Scénario} & \textbf{Ce qui est vérifié} \\
  \midrule
  1 & Groupe de 5 étudiants       & Moyenne, médiane et résultat complet \\
  2 & Notes toutes identiques     & Moyenne = médiane = note commune \\
  3 & Grand groupe de 10 (pair)   & Médiane = moyenne des deux centraux \\
  4 & Groupe d'un seul étudiant   & Cas limite : moyenne = médiane = note \\
  5 & Identifiants de l'enseignant & Email valide accepté, invalide rejeté \\
  \bottomrule
\end{tabular}
\end{center}

\subsection{Extrait --- Scénario 1}

\begin{lstlisting}[style=java, caption={\texttt{ScenarioGroupeTPTest.java} --- scénario 1}]
@Test
@DisplayName("Le resultat stocke contient toutes les informations requises")
void givenGroupeDeCinqEtudiants_whenGenerateResultat_thenContainsAllInfos() {
    // Given
    UtilisateurTab enseignant = new UtilisateurTab(
            "Alain", "Dupont", "alain.dupont@iut.fr", DATE_EXAMEN);
    CalculTab calcul = new CalculTab();
    for (int note : new int[]{10, 12, 14, 16, 18}) calcul.ajouterNote(note);
    // When
    String resultat = calcul.toResultatString(enseignant);
    // Then
    assertThat(resultat)
            .contains("Alain").contains("Dupont")
            .contains("alain.dupont@iut.fr").contains("2026-05-13")
            .contains("5").contains("14.00");
}
\end{lstlisting}

\subsection{État rouge --- avant \texttt{toResultatString()}}

\etatrouge{2 tests de scénario échouent : \texttt{toResultatString()} retourne
une chaîne vide.}

\screenshot{AE_04_scenarios_rouge.png}{%
  \texttt{mvn test -Dtest=ScenarioGroupeTPTest} --- avant \texttt{toResultatString()}}

\subsection{Implémentation}

\begin{lstlisting}[style=java, caption={\texttt{CalculTab.java} --- toResultatString()}]
public String toResultatString(UtilisateurTab utilisateur) {
    return new StringBuilder()
            .append(utilisateur)      // prenom, nom, email, date
            .append(", ").append(getTaille())
            .append(", ").append(String.format(Locale.ROOT, "%.2f",
                    calculerMoyenne()))
            .append(", ").append(String.format(Locale.ROOT, "%.2f",
                    calculerMediane()))
            .toString();
}
\end{lstlisting}

\subsection{État vert --- scénarios complets}

\etatvert{8 tests de scénario passent --- 5 scénarios couverts.}

\screenshot{AE_04_scenarios_vert.png}{%
  \texttt{mvn test -Dtest=ScenarioGroupeTPTest} --- scénarios au vert}

\subsection{Bilan global}

\etatvert{25 tests --- 0 échec --- BUILD SUCCESS}

\screenshot{AE_05_final_vert.png}{%
  \texttt{mvn test} --- exécution complète : 25 tests au vert}

% =====================================================================
\section{Extension du projet}
% =====================================================================

Cette partie complète l'application avec la validation des plages de notes
et la persistance du résultat dans un fichier.

% ----------------------------
\subsection{Validation des notes (\texttt{ajouterNote})}
% ----------------------------

\subsubsection*{Objectif}

Une note doit être comprise entre 0 et 20 inclus. Toute valeur hors plage
doit lever une \texttt{IllegalArgumentException}.

\subsubsection*{Tests prévus}

\begin{lstlisting}[style=java, caption={\texttt{CalculTabEdgeCasesTest.java} --- extrait attendu}]
@Test
void ajouterNote_throwsException_whenNoteAbove20() {
    assertThatThrownBy(() -> calcul.ajouterNote(25))
            .isInstanceOf(IllegalArgumentException.class);
}

@Test
void ajouterNote_throwsException_whenNoteNegative() {
    assertThatThrownBy(() -> calcul.ajouterNote(-1))
            .isInstanceOf(IllegalArgumentException.class);
}

@Test
void ajouterNote_succeeds_whenNoteIsZero() {
    calcul.ajouterNote(0);
    assertThat(calcul.getTaille()).isEqualTo(1);
}
\end{lstlisting}

\subsubsection*{État rouge}

\etatrouge{\texttt{ajouterNote(25)} s'exécute sans erreur --- aucune
validation n'est présente.}

\screenshot{LC_01_validation_rouge.png}{%
  Tests de validation des notes en échec}

\subsubsection*{Implémentation attendue}

\begin{lstlisting}[style=java, caption={\texttt{CalculTab.java} --- garde dans ajouterNote()}]
public void ajouterNote(int note) {
    if (note < 0 || note > 20)
        throw new IllegalArgumentException(
                "Note invalide : " + note + " (doit etre entre 0 et 20)");
    notes.add(note);
}
\end{lstlisting}

\subsubsection*{État vert}

\etatvert{Tests de validation des notes passent.}

\screenshot{LC_01_validation_vert.png}{%
  Tests de validation des notes au vert}

% ----------------------------
\subsection{Stockage du résultat (\texttt{ResultatStockage})}
% ----------------------------

\subsubsection*{Objectif}

La classe \texttt{ResultatStockage} écrit la chaîne retournée par
\texttt{toResultatString()} dans un fichier texte. Les tests vérifient
que le fichier est créé et contient le contenu exact.

\subsubsection*{État rouge}

\etatrouge{Erreur de compilation : la classe \texttt{ResultatStockage}
n'existe pas encore.}

\screenshot{LC_02_stockage_rouge.png}{%
  Erreur de compilation avant la création de \texttt{ResultatStockage}}

\subsubsection*{Interface attendue}

\begin{lstlisting}[style=java, caption={\texttt{ResultatStockage.java} --- interface attendue}]
public class ResultatStockage {
    /**
     * Ecrit le contenu dans le fichier au chemin indique.
     * @param contenu  chaine retournee par toResultatString()
     * @param chemin   chemin du fichier cible
     */
    public void sauvegarder(String contenu, String chemin) throws IOException {
        // a implementer
    }
}
\end{lstlisting}

\subsubsection*{État vert}

\etatvert{Tests de stockage passent : fichier créé avec le bon contenu.}

\screenshot{LC_02_stockage_vert.png}{%
  Tests de stockage au vert}

% ----------------------------
\subsection{Résultat final du projet}
% ----------------------------

\screenshot{LC_03_final_vert.png}{%
  \texttt{mvn test} --- tous les tests passent après la partie LC}

% =====================================================================
\section{Conclusion}
% =====================================================================

Ce projet nous a permis de mettre en pratique la démarche \textbf{TDD}
sur un cas concret. En écrivant les tests avant le code, nous avons été
contraints de définir précisément le comportement attendu avant toute
implémentation.

Les points clés retenus :

\begin{itemize}
  \item Le cycle \textbf{Rouge $\rightarrow$ Vert} pousse à écrire un code
        minimal et ciblé, sans surconception.
  \item Les \textbf{scénarios \texttt{@Nested}} donnent une vision claire
        des cas d'usage réels plutôt que de simples tests unitaires isolés.
  \item \textbf{AssertJ} produit des messages d'erreur explicites qui
        accélèrent le diagnostic en cas d'échec.
  \item La \textbf{séparation des classes} (\texttt{UtilisateurTab},
        \texttt{CalculTab}, \texttt{Principale}) facilite les tests unitaires
        en évitant les dépendances inutiles.
  \item Les \textbf{commits atomiques} permettent de retrouver l'état du
        projet à chaque étape et de visualiser la progression via l'historique
        Git.
\end{itemize}

Au final, \textbf{25 tests} couvrent l'ensemble des fonctionnalités
principales. La couverture dépasse largement le seuil de 50\,\% demandé.
La partie LC vient compléter ce bilan avec la validation des plages de
notes et la persistance du résultat.

\end{document}
